\documentclass[11pt]{article}
\usepackage[margin=1in]{geometry}
\usepackage[T1]{fontenc}
\usepackage[utf8]{inputenc}
\usepackage{lmodern}
\usepackage{microtype}
\usepackage{xcolor}
\usepackage{amsmath,amssymb,amsthm}
\usepackage{enumitem}
\usepackage{hyperref}
\usepackage{setspace}
\usepackage{titlesec}
\usepackage{parskip}
\usepackage{booktabs}

\hypersetup{colorlinks=true,linkcolor=black,urlcolor=blue,citecolor=black}
\setstretch{1.08}
\setlist[itemize]{topsep=4pt,itemsep=2pt,leftmargin=1.4em}
\titleformat{\section}{\large\bfseries}{\thesection.}{0.6em}{}
\titleformat{\subsection}{\normalsize\bfseries}{\thesubsection.}{0.5em}{}

\newtheorem{theorem}{Theorem}
\newtheorem{proposition}[theorem]{Proposition}
\newtheorem{lemma}[theorem]{Lemma}
\newtheorem{corollary}[theorem]{Corollary}
\newtheorem{definition}[theorem]{Definition}
\newtheorem{remark}[theorem]{Remark}
\newtheorem{observation}[theorem]{Observation}

\newcommand{\ALCIRCC}[1]{\ensuremath{\mathcal{ALCI}_{\mathit{RCC#1}}}}
\newcommand{\LRCC}[1]{\ensuremath{L_{\mathit{RCC#1}}}}
\newcommand{\ALC}{\ensuremath{\mathcal{ALC}}}
\newcommand{\ALCI}{\ensuremath{\mathcal{ALCI}}}
\newcommand{\NR}{\ensuremath{\mathcal{N}_\mathcal{R}}}
\newcommand{\NC}{\ensuremath{\mathcal{N}_\mathcal{C}}}
\newcommand{\reals}{\mathbb{R}}
\newcommand{\naturals}{\mathbb{N}}
\newcommand{\EXPTIME}{\textsc{ExpTime}}
\newcommand{\PSPACE}{\textsc{PSpace}}
\newcommand{\comp}{\mathsf{comp}}
\newcommand{\tp}{\mathsf{tp}}
\newcommand{\RS}{\mathit{RS}}
\newcommand{\RSE}{\mathit{RS}^{\exists}}

\title{\bfseries The $\LRCC{8}$ Undecidability Proof Does Not Transfer to $\ALCIRCC{8}$:\\
  Concrete vs.\ Abstract Semantics and the Coincidence Obstruction}

\author{Claude (Anthropic) and Michael Wessel\thanks{%
  Michael Wessel introduced the $\ALCIRCC{}$ family and identified the
  coincidence obstruction in~\cite{Wessel2003,Wessel2001}.
  Claude performed the comparative analysis and wrote this report.
  Contact: miacwess@gmail.com.}}

\date{April 2026}

\begin{document}
\maketitle

\begin{abstract}
Lutz and Wolter~\cite{LutzWolter2006} proved that the modal logic $\LRCC{8}$
is undecidable by reducing the $\naturals\times\naturals$ domino tiling problem
to satisfiability over topological models.  A natural question is whether this
undecidability transfers to $\ALCIRCC{8}$, the description logic with RCC8 roles
and composition-table semantics introduced by Wessel~\cite{Wessel2003}.

We show that it does not.  The reduction relies on the geometry of regular closed
sets in~$\reals^2$ to enforce a \emph{grid coincidence condition}
($R_X \circ R_Y = R_Y \circ R_X$) that the abstract composition-table semantics
of $\ALCIRCC{8}$ cannot guarantee.  This coincidence obstruction was first
identified by Wessel in~2002, who constructed explicit RCC8 grid models but
showed that no concept term can force the coincidence at the level of abstract
models.

The analysis reveals that Wessel's 2002/2003 technical report anticipated the
key technical ingredients of the Lutz--Wolter proof---the even-odd chain
construction, the grid encoding via TPP/NTPP, and the identification of the
coincidence gap as the decisive obstacle---three to four years before their
LMCS publication.  The decidability of $\ALCIRCC{8}$ remains open.
\end{abstract}

%% ====================================================================
\section{Introduction}
%% ====================================================================

The description logics $\ALCIRCC{5}$ and $\ALCIRCC{8}$ were introduced by
Wessel~\cite{Wessel2001,Wessel2003} as extensions of $\ALCI$ with spatially
interpreted roles.  In these logics, every pair of domain elements is related
by exactly one RCC base relation, the role box is fixed by the RCC composition
table, and the identity relation is the only reflexive role (strong EQ
semantics).  Wessel proved $\ALCIRCC{1}$, $\ALCIRCC{2}$, and $\ALCIRCC{3}$
decidable, and established lower complexity bounds for the remaining two:
$\PSPACE$-hard for $\ALCIRCC{5}$ and $\EXPTIME$-hard for
$\ALCIRCC{8}$~\cite{Wessel2003}.  The decidability of $\ALCIRCC{5}$ and
$\ALCIRCC{8}$ was left as an open problem and remains open to this day.

Independently, Lutz and Wolter~\cite{LutzWolter2006} studied \emph{modal
logics of topological relations}, denoted $\LRCC{8}$, $\LRCC{5}$, etc.  These
logics have the same \emph{syntax} as the $\ALCIRCC{}$ family but are
interpreted over \emph{topological models}: the domain elements are regular
closed subsets of a topological space (typically~$\reals^n$), and the RCC
relations have their standard topological definitions.  They proved $\LRCC{8}$
undecidable by reducing the $\naturals\times\naturals$ domino tiling
problem~\cite{LutzWolter2006}.

The syntactic coincidence naturally raises the question:

\begin{quote}
\emph{Does the undecidability of $\LRCC{8}$ imply the undecidability
of~$\ALCIRCC{8}$?}
\end{quote}

We show that the answer is \textbf{no}.  The core of the argument is the
\emph{coincidence obstruction}: the Lutz--Wolter reduction requires that
certain pairs of regions coincide (to close the grid), and this coincidence is
guaranteed by the geometry of~$\reals^2$ but not by the composition table
alone.  Wessel identified exactly this obstruction in~\cite{Wessel2003},
Section~4.4.3, and gave explicit constructions showing that the grid can be
built as an RCC8 network but cannot be \emph{enforced} by a concept term in the
abstract semantics.

The purpose of this note is threefold:
\begin{enumerate}
  \item To make the non-transfer argument precise (Section~\ref{sec:nontransfer}).
  \item To document the technical priority of the constructions in
    Wessel~\cite{Wessel2003} relative to Lutz--Wolter~\cite{LutzWolter2006}
    (Section~\ref{sec:priority}).
  \item To discuss what the comparison tells us about the decidability status of
    $\ALCIRCC{8}$ (Section~\ref{sec:implications}).
\end{enumerate}

%% ====================================================================
\section{Preliminaries}
\label{sec:prelim}
%% ====================================================================

\subsection{RCC8 base relations and composition table}

RCC8 distinguishes eight base relations between regions:
\[
\NR_8 = \{DC,\, EC,\, PO,\, EQ,\, TPP,\, NTPP,\, TPPI,\, NTPPI\}.
\]
$DC$ (disconnected) and $EC$ (externally connected) refine $DR$; $TPP$
(tangential proper part) and $NTPP$ (non-tangential proper part) refine $PP$.
The composition table $\comp(R,S) \subseteq \NR_8$ is defined for all pairs
$R,S \in \NR_8$ and captures which base relations are possible for a pair
$(a,c)$ given that $R(a,b)$ and $S(b,c)$ hold.

A key property of the composition table is:
\begin{align}
  TPPI \circ TPPI &\;\sqsubseteq\; TPPI \sqcup NTPPI, \label{eq:tppi-comp}\\
  TPPI \circ NTPPI &\;\sqsubseteq\; NTPPI, \label{eq:tppi-ntppi}\\
  NTPPI \circ TPPI &\;\sqsubseteq\; NTPPI, \label{eq:ntppi-tppi}\\
  NTPPI \circ NTPPI &\;\sqsubseteq\; NTPPI. \label{eq:ntppi-comp}
\end{align}
These entries mean that the composition of the \emph{transitive orbit} of $TPPI$
with itself is contained in the transitive orbit plus $NTPPI$: in other words,
$NTPPI$ is the ``non-immediate'' part of the transitive closure of $TPPI$.  This
is what allows the distinction between \emph{direct} and \emph{indirect}
containment---a distinction unavailable in RCC5, where $PP = TPP \cup NTPP$ is
a single relation.

\subsection{$\ALCIRCC{8}$: abstract composition-table semantics}

An $\ALCIRCC{8}$ interpretation is a pair $\mathcal{I} = (\Delta^{\mathcal{I}},
\cdot^{\mathcal{I}})$ where:
\begin{itemize}
  \item $\Delta^{\mathcal{I}}$ is a non-empty set (the domain);
  \item every pair $(d,e)$ with $d \neq e$ is assigned exactly one base
    relation $R \in \NR_8 \setminus \{EQ\}$, and $EQ^{\mathcal{I}} =
    \mathrm{Id}(\Delta^{\mathcal{I}})$
    (strong EQ semantics);
  \item the role box (composition-table constraints) is satisfied:
    for all $R,S \in \NR_8$, $R^{\mathcal{I}} \circ S^{\mathcal{I}} \subseteq
    \bigcup_{T \in \comp(R,S)} T^{\mathcal{I}}$;
  \item concepts are interpreted as in $\ALCI$: $(\exists R.C)^{\mathcal{I}} =
    \{d \mid \exists e.\; \langle d,e\rangle \in R^{\mathcal{I}} \wedge e \in
    C^{\mathcal{I}}\}$, etc.
\end{itemize}
The key point: an $\ALCIRCC{8}$ model is an \emph{abstract complete graph}
with RCC8 edge-coloring satisfying the composition table.  No topological
realization is required.

\subsection{$\LRCC{8}$: topological semantics}

An $\LRCC{8}$ model~\cite{LutzWolter2006} interprets formulas over regular
closed subsets of a topological space (typically~$\reals^n$).  The RCC8
relations are defined topologically:
\begin{itemize}
  \item $DC(a,b)$ iff $a \cap b = \emptyset$;
  \item $EC(a,b)$ iff $a \cap b \neq \emptyset$ but $a^\circ \cap b^\circ = \emptyset$;
  \item $PO(a,b)$ iff $a^\circ \cap b^\circ \neq \emptyset$, $a \not\subseteq b$,
    $b \not\subseteq a$;
  \item $TPP(a,b)$ iff $a \subsetneq b$ and $a \cap \partial b \neq \emptyset$;
  \item $NTPP(a,b)$ iff $a \subsetneq b$ and $a \cap \partial b = \emptyset$;
\end{itemize}
and symmetrically for $TPPI$, $NTPPI$, and $EQ$.

Every topological model is an $\ALCIRCC{8}$ model (the composition table holds
in any topological space), but not every $\ALCIRCC{8}$ model is topologically
realizable.  Formally:

\begin{observation}\label{obs:inclusion}
  For any concept/formula $\varphi$:
  \[
  \varphi \text{ is } \LRCC{8}\text{-satisfiable}
  \;\;\Longrightarrow\;\;
  \varphi \text{ is } \ALCIRCC{8}\text{-satisfiable}.
  \]
  The converse does not hold in general.
\end{observation}

%% ====================================================================
\section{The even-odd chain: functional successors via TPP/NTPP}
\label{sec:chain}
%% ====================================================================

Both Wessel~\cite{Wessel2003} and Lutz--Wolter~\cite{LutzWolter2006} exploit
the TPP/NTPP distinction to create \emph{functional successor chains} within
RCC8 models.  The construction deserves careful comparison.

\subsection{Wessel's construction (2002/2003)}

In~\cite{Wessel2003}, Section~4.4.2, Wessel defines the concept
\texttt{infinite\_even\_odd\_chain}:
\[
  \mathit{even} \;\sqcap\; (\mathit{even} \to \forall TPPI.\mathit{odd})
  \;\sqcap\; \forall R_*.\bigl((\mathit{even} \leftrightarrow \neg\mathit{odd})
  \;\sqcap\; (\mathit{even} \to \forall TPPI.\mathit{odd})
  \;\sqcap\; (\mathit{odd} \to \forall TPPI.\mathit{even})\bigr)
  \;\sqcap\; \cdots
\]
(simplified; see~\cite{Wessel2003}, p.~39 for the full definition).
Each node alternates between $\mathit{even}$ and $\mathit{odd}$.
Since~\eqref{eq:tppi-comp}--\eqref{eq:ntppi-comp} imply
$(TPPI^{\mathcal{I}})^+ - TPPI^{\mathcal{I}} \subseteq NTPPI^{\mathcal{I}}$,
a node can distinguish its \emph{direct} $TPPI$-successor from all
\emph{indirect} ones by parity.  This gives a \emph{functional} successor
relation: each node has exactly one $TPPI$-successor of the opposite parity.

Wessel then uses this chain to encode a \textbf{binary $n$-bit counter},
showing that $\ALCIRCC{8}$ models can be forced to have paths of length
$\geq 2^n$, thereby proving:

\begin{theorem}[Wessel~\cite{Wessel2003}, Section~4.4.2]
  Concept satisfiability in $\ALCIRCC{8}$ is $\EXPTIME$-hard.
\end{theorem}

\subsection{Lutz--Wolter's construction (2006)}

In~\cite{LutzWolter2006}, the same idea appears as follows.  An infinite
sequence of regions $r_1, r_2, r_3, \ldots$ is ordered by $\mathit{ntpp}$.
Regions are marked alternately with $a \wedge b$ and $a \wedge \neg b$,
connected via $\mathit{tpp}$.  The \emph{next} operator is defined as:
\[
  \mathsf{3}^+ \varphi \;=\; \langle tppi\rangle(a \wedge \neg b \wedge
  \langle tppi\rangle(a \wedge b \wedge \varphi))
\]
which skips one intermediate region to reach the next grid position.  This is
the same functional-successor construction, using $tpp/ntpp$ in place of
$TPP/NTPP$ and $a \wedge b$ / $a \wedge \neg b$ in place of
$\mathit{even}$/$\mathit{odd}$.

Lutz and Wolter credit this technique to Marx and Reynolds~\cite{MarxReynolds1999}
and Rybina and Zakharyaschev~\cite{RybinaZakharyaschev2001}.  They do not
reference Wessel's even-odd chain construction~\cite{Wessel2003}, despite
citing Wessel's earlier workshop paper~\cite{Wessel2001}.

\begin{remark}
  Wessel's even-odd chain appears in his DL Workshop 2001 paper~\cite{Wessel2001}
  and is developed more fully in the 2002/2003 technical
  report~\cite{Wessel2003}.  The technique of distinguishing $TPP$ from $NTPP$
  to obtain functional successors in an RCC8 logic is independently present in
  Marx--Reynolds~\cite{MarxReynolds1999} for temporal logics.  The
  application to RCC8 description logics---including the binary counter and the
  $\EXPTIME$-hardness proof---is Wessel's contribution.
\end{remark}

%% ====================================================================
\section{The grid construction and the coincidence obstruction}
\label{sec:grid}
%% ====================================================================

Both Wessel and Lutz--Wolter attempt to encode a two-dimensional grid
$\naturals \times \naturals$ in RCC8 models, in order to reduce the domino
tiling problem.  The grid requires:
\begin{enumerate}[label=(G\arabic*)]
  \item\label{g:nodes} Nodes represent grid positions.
  \item\label{g:func} Each node has a unique horizontal successor ($R_X$) and
    a unique vertical successor ($R_Y$), both extractable from $TPP$.
  \item\label{g:coinc} \textbf{Coincidence condition:}
    $R_X \circ R_Y = R_Y \circ R_X$.
    That is, the east-of-north equals the north-of-east.
\end{enumerate}

\subsection{Wessel's grid construction and the obstruction (2002/2003)}

In~\cite{Wessel2003}, Section~4.4.3, Wessel:

\begin{enumerate}
  \item Defines the domino system and the requirements for a grid concept
    $C_{\mathit{grid}}$ (p.~41).

  \item Constructs explicit RCC8 networks isomorphic to
    $n \times n$ grids for arbitrary~$n$ (Figure~10, p.~42).
    The construction uses $TPP$ for horizontal and vertical successors,
    $PO$ for within-row relationships, and $NTPP$ for indirect containment.
    These networks are composition-consistent.

  \item Identifies the \textbf{coincidence obstruction} (Figure~11, p.~43):
    when a node~$x$ has two $TPP$-successors~$y$ (horizontal) and another
    (vertical), and~$y$ has a $TPP$-successor~$z$ (vertical), and the
    vertical path from~$x$ leads to $z'$ (horizontal), then enforcing
    $z = z'$ requires $EQ(z,z')$.  But the composition table allows
    \[
    \{PO, EQ, TPP, NTPP, TPPI, NTPPI\}
    \]
    as possible relations between~$z$ and~$z'$---there is no way to
    force $EQ$ using concept terms alone.

  \item Concludes (p.~43--44):
    \begin{quote}
      ``Even though these infinite models do exist, we have found no way to
      \emph{enforce} the depicted model(s) by means of a concept term.''
    \end{quote}

  \item Observes that adding a hybrid logic binding operator $\downarrow$
    would immediately yield undecidability, confirming that the coincidence
    condition is the \emph{precise boundary}.
\end{enumerate}

\subsection{Lutz--Wolter's grid construction (2006)}

In~\cite{LutzWolter2006}, Section~8, Lutz and Wolter solve the coincidence
problem---but only in the \emph{topological} setting.  They introduce
additional ``$c$-regions'' connected via $tpp$ to implement a ``go right''
function in the grid.  The grid coincidence~\ref{g:coinc} is enforced
because:
\begin{itemize}
  \item In a topological space, $TPP(a,b)$ means $a \subsetneq b$ with
    $a \cap \partial b \neq \emptyset$.  The containment structure of regions
    in~$\reals^2$ (or~$\reals^n$ for $n \geq 1$) is rigid enough that the
    ``go right then go up'' and ``go up then go right'' paths are forced
    to reach the same region.
  \item This rigidity is a consequence of the topological structure, not of the
    composition table.  The composition table is a \emph{necessary} condition
    for topological realizability but not a \emph{sufficient} one: there exist
    infinite composition-consistent RCC8 networks that are not realizable in
    any topological space.
\end{itemize}

The result is:

\begin{theorem}[Lutz--Wolter~\cite{LutzWolter2006}, Theorem~8.2]
  Satisfiability of $\LRCC{8}$ formulas (over topological models in~$\reals^n$
  for any $n \geq 1$) is undecidable.
\end{theorem}

%% ====================================================================
\section{Why the undecidability does not transfer}
\label{sec:nontransfer}
%% ====================================================================

We now make the non-transfer argument precise.

\begin{proposition}\label{prop:nontransfer}
  The undecidability of $\LRCC{8}$-satisfiability does not imply the
  undecidability of $\ALCIRCC{8}$-satisfiability.  Specifically, the
  Lutz--Wolter reduction from the $\naturals \times \naturals$ domino tiling
  problem cannot be lifted from topological to abstract models.
\end{proposition}

\begin{proof}
  Let $\mathcal{D} = (\mathcal{D}, \mathcal{H}, \mathcal{V})$ be a domino
  system, and let $C_\mathcal{D}$ be the Lutz--Wolter formula encoding
  $\mathcal{D}$.  The reduction requires:
  \[
  \mathcal{D} \text{ tiles } \naturals \times \naturals
  \quad\Longleftrightarrow\quad
  C_\mathcal{D} \text{ is satisfiable}
  \]
  in the relevant model class.

  \textbf{Soundness ($\Rightarrow$):}
  If $\mathcal{D}$ tiles, then $C_\mathcal{D}$ has a topological model (the
  grid of regions in~$\reals^2$).  By Observation~\ref{obs:inclusion}, it
  also has an abstract model.  So $C_\mathcal{D}$ is $\ALCIRCC{8}$-satisfiable.
  This direction transfers.

  \textbf{Completeness ($\Leftarrow$):}
  If $\mathcal{D}$ does \emph{not} tile, we need $C_\mathcal{D}$ to be
  unsatisfiable in the relevant model class.  For $\LRCC{8}$, every model is
  topological, and the topological rigidity forces the grid coincidence
  condition~\ref{g:coinc}.  Therefore every $\LRCC{8}$-model of
  $C_\mathcal{D}$ is (isomorphic to) a valid tiling, so
  $\LRCC{8}$-unsatisfiability follows.

  For $\ALCIRCC{8}$, the situation is different.  An abstract model of
  $C_\mathcal{D}$ need not satisfy~\ref{g:coinc}: the ``east-of-north''
  and ``north-of-east'' elements $z$ and $z'$ may be distinct, with
  $R(z,z') \in \{PO, TPP, NTPP, TPPI, NTPPI\}$ (any relation permitted
  by the composition table).  Such a model satisfies the concept term
  $C_\mathcal{D}$ without encoding a valid tiling---it is a ``spurious''
  model where the two-dimensional grid structure unravels.

  Therefore, $C_\mathcal{D}$ being $\ALCIRCC{8}$-satisfiable does not
  imply that $\mathcal{D}$ tiles, and the completeness direction fails.
  \qedhere
\end{proof}

\begin{remark}
  The failure is not merely a gap in the proof that might be closed by a
  cleverer encoding.  It reflects a \emph{structural} difference between the
  two semantics: topological models have geometric rigidity that abstract
  models lack.  Any domino reduction to $\ALCIRCC{8}$ would need to enforce
  grid coincidence using only the composition table---but as Wessel
  showed~\cite{Wessel2003}, the table allows too many options at the
  coincidence point (Figure~11).
\end{remark}

\begin{remark}\label{rem:rcc5-parallel}
  The situation is structurally parallel to RCC5.  Lutz and
  Wolter~\cite{LutzWolter2006} prove $\LRCC{5}(\RSE)$ undecidable (via a
  reduction from $S5^3$ satisfiability), where $\RSE$ denotes region
  structures with supremum closure.  This does not transfer to
  $\ALCIRCC{5} = \LRCC{5}(\RS)$ because abstract $\ALCIRCC{5}$ models need
  not have suprema.  Lutz and Wolter explicitly acknowledge this:
  $\LRCC{5}(\RS)$ is ``one of the most intriguing open problems'' suggested
  by their work~\cite{LutzWolter2006}.  Both non-transfer arguments have
  the same shape: the concrete semantics (topological for RCC8, supremum
  closure for RCC5) provides constraints beyond the composition table that
  the abstract semantics lacks.
\end{remark}

%% ====================================================================
\section{Could the reduction be adapted?}
\label{sec:adapt}
%% ====================================================================

One might ask: can the Lutz--Wolter encoding be modified to work in the
abstract setting?  We consider three natural approaches.

\subsection{Using nominals to force coincidence}

As Wessel observes~\cite{Wessel2003}, $\ALCIRCC{8}$ can express nominals:
the concept $\boxed{i} \sqcap \forall R_D.\neg\boxed{i}$ (where
$R_D = \NR \setminus \{EQ\}$) forces $\boxed{i}^{\mathcal{I}}$ to be a
singleton.  One could attempt to name the ``diagonal'' node $z$ and $z'$ with
the same nominal, forcing $z = z'$.

But this works only for \emph{finitely many} grid positions.  The domino
problem requires an \emph{infinite} grid, and a concept term contains only
finitely many nominals.  As Wessel notes (p.~43--44): ``this identification
[\ldots] works only for a finite number of successors, since we can only
exploit a finite number of nominals.''

Adding a hybrid logic binding operator $\downarrow$ (which can
introduce fresh nominals dynamically) would solve the problem---but also
immediately yield undecidability, as this goes beyond
$\ALCIRCC{8}$~\cite{Wessel2003}.

\subsection{Exploiting patchwork to close the gap}

The RCC8 composition table satisfies a \emph{patchwork property}: every
path-consistent atomic constraint network is globally
consistent~\cite{RenzNebel1999}.  Could this be used to argue that all
abstract models satisfying $C_\mathcal{D}$ must be grid-like?

No.  Patchwork guarantees that a \emph{consistent atomic network can be
extended}, but it does not force two nodes to be \emph{equal}.  The
coincidence condition requires $z = z'$ (i.e., $EQ(z,z')$), not just
consistency of the relation between them.  Patchwork is a tool for
\emph{soundness} (building models), not for \emph{rigidity} (constraining
models).

\subsection{Using the alternating-type trick}

In~\cite{OurMainPaper}, we observe that $\ALCIRCC{8}$ can distinguish a
role from its transitive orbit: $\forall TPP.C$ affects only direct
$TPP$-successors, while $\forall NTPP.C$ affects indirect ones.  This
allows an \emph{alternating-type trick}: a concept $\forall NTPP.D$ broadcasts
$D$ to all indirect successors, and $D$ can constrain the types of nodes
reached two or more steps along the $TPPI$-chain.

This gives $\ALCIRCC{8}$ a limited form of ``global constraint propagation.''
But propagation is \emph{monotonic} (constraints accumulate, they cannot be
consumed), which is insufficient for the unbounded synchronization required by
general domino tiling.  The alternating-type trick provides a 1-dimensional
ladder but not a 2-dimensional grid.

\subsection{A concrete $TPPI/NTPPI$ grid attempt}
\label{sec:concrete-tppi-attempt}

To make the obstruction completely explicit, we write out a concrete
$\ALCIRCC{8}$ concept that attempts to encode a 2D grid using the
$TPPI/NTPPI$ distinction.  The construction is the natural RCC8 lift
of the failed RCC5 domino attempt in our main paper~\cite{OurMainPaper}:
there, $\comp(PPI,PPI) = \{PPI\}$ makes $PPI$-chains transitive and
the ancestral isolation universal propagates to every deeper
descendant.  The RCC8 refinement $\comp(TPPI,TPPI) = \{TPPI,NTPPI\}$
is non-transitive, which suggests the collapse can be avoided.  We
show that while this is locally true, the \emph{coincidence
obstruction} of~\S\ref{sec:grid} reappears at the grid level.

\medskip\noindent\textbf{Schema.}  Let $\mathsf{Corner}$ be a
\emph{local} isolation marker placed directly on the colored
grandchildren (rather than as an ancestral universal, so it cannot
propagate transitively through recursion):
\[
  \mathsf{Corner} \;=\;
  \bigwedge_{R \,\in\, \{DC,EC,PO,TPP,NTPP,TPPI,NTPPI\}}
    \forall R.(\neg C \sqcap \neg D).
\]
Define
\begin{align*}
  X \;=\; &\exists TPPI.\exists TPPI.(C \sqcap \mathsf{Corner})
    \;\sqcap\; \exists TPPI.\exists TPPI.(D \sqcap \mathsf{Corner}) \\
  &\sqcap\; \forall TPPI.(\neg C \sqcap \neg D)
    \;\sqcap\; \forall TPPI.X.
\end{align*}
The conjunct $\forall TPPI.(\neg C \sqcap \neg D)$ forces the corner
relation $x \to z$ to be $NTPPI$ (not $TPPI$), via
$\comp(TPPI,TPPI) = \{TPPI,NTPPI\}$.

\medskip\noindent\textbf{Level 1 is satisfiable, with direct/indirect
successors genuinely distinguished.}  Without recursion, $x$ has two
$TPPI$-children $y_1, y_2$ and a single $NTPPI$-grandchild $z$
reached along both $TPPI$-$TPPI$ chains.  The RCC5 EQ-collapse
argument carries over: $z_1 \in C \sqcap \mathsf{Corner}$ and
$z_2 \in D \sqcap \mathsf{Corner}$ are mutual non-$EQ$ neighbors only
if $\mathsf{Corner}$ is violated, so under strong-$EQ$ semantics
$z_1 = z_2 = z$.  Crucially, $z$ is $NTPPI$-related to $x$ while
$y_1, y_2$ are $TPPI$-related---the direct vs.~indirect distinction
survives.  This is a genuine refinement over the RCC5 attempt
in~\cite{OurMainPaper}, where $PP/PPI$-chains cannot express such a
distinction.

\medskip\noindent\textbf{Level 2 collapses via the coincidence
obstruction.}  Adding $\forall TPPI.X$: the $TPPI$-child $y$ satisfies
$X$, producing its own corner $w \in C \sqcap D \sqcap \mathsf{Corner}$
at $NTPPI$-distance from $y$, and hence at $NTPPI$-distance from $x$
via $\comp(TPPI,NTPPI) = \{NTPPI\}$.  Now $z$ (the corner of $x$)
and $w$ (the corner of $y$) are distinct \emph{intended} grid
positions, but as elements of the model they are related by some RCC8
base relation.  Since $z$ carries $\mathsf{Corner}$, every non-$EQ$
neighbor of $z$ lies in $\neg C \sqcap \neg D$.  But $w \in C \sqcap D$.
Therefore $z = w$ under strong-$EQ$ semantics.  Iterating the
recursion, \textbf{all corners of the intended grid coincide at a
single point}---the grid collapses to 1D.

\medskip\noindent\textbf{Diagnosis.}  The RCC5 collapse is a
\emph{transitive-propagation} failure: the grandparent universal fires
on great-grandchildren because $\comp(PPI,PPI) = \{PPI\}$.  Moving
the isolation into a local $\mathsf{Corner}$ marker and using the
non-transitive composition $\comp(TPPI,TPPI) = \{TPPI,NTPPI\}$
eliminates that pathway.  But a different obstruction takes over: the
composition table permits any of
$\{PO, EQ, TPP, NTPP, TPPI, NTPPI\}$ between two corners at
distinct intended grid coordinates, and nothing in the concept
language rules out $EQ$.  The local $\mathsf{Corner}$ marker then
\emph{forces} the $EQ$-merge.  This is exactly the composition-table
gap that Wessel's Figure~11 identifies in~\cite{Wessel2003}
(cf.~\S\ref{sec:grid}): abstract RCC8 does not carry the rigidity
needed to keep distinct grid cells distinct.

\begin{remark}\label{rem:concrete-attempt-conclusion}
  This concrete attempt sharpens the non-transfer story.  The
  $TPPI/NTPPI$ distinction does provide something the RCC5 case
  lacks---a local direct/indirect separator at level~2---and it does
  avoid the specific transitive-propagation failure mode of the RCC5
  attempt.  But it does \emph{not} close the 2D coincidence gap.
  Lutz--Wolter's $\LRCC{8}$ reduction resolves the gap topologically:
  geometric boundary inclusions force north-of-east to coincide with
  east-of-north \emph{without} causing distant corners to merge, because
  cells are 2D regions with distinct interiors~\cite{LutzWolter2006}.
  No such rigidity is available in the abstract setting, consistent
  with $\ALCIRCC{8}$ decidability remaining open.
\end{remark}

%% ====================================================================
\section{Technical priority and citation analysis}
\label{sec:priority}
%% ====================================================================

Lutz and Wolter~\cite{LutzWolter2006} cite one paper by
Wessel---\cite{Wessel2001}, the DL Workshop 2001 paper.  They do \emph{not}
cite the 2002/2003 technical report~\cite{Wessel2003}, which contains the
more developed results.  Their characterization of Wessel's work is:

\begin{quote}
  ``We should note that modal logics based on the Egenhofer-Franzosa
  relations have been suggested in an early paper by Cohn [Coh93] and
  further considered in [Wes01].  However, it proved difficult to analyze
  the expressive power and computational behavior of such logics: despite
  several efforts, to the best of our knowledge no results have been
  obtained so far.''~\cite{LutzWolter2006}
\end{quote}

This characterization is incomplete in several respects:

\begin{enumerate}
  \item \textbf{``No results'':}
    Even the cited paper~\cite{Wessel2001} contains undecidability results
    for $\mathcal{ALC}_{\mathcal{RA}^\ominus}$ (ALC with arbitrary role
    axioms, via PCP reduction).  The uncited technical
    report~\cite{Wessel2003} goes significantly further: it proves
    decidability of $\ALCIRCC{1}$, $\ALCIRCC{2}$, and $\ALCIRCC{3}$;
    establishes $\PSPACE$-hardness of $\ALCIRCC{5}$ and $\EXPTIME$-hardness
    of $\ALCIRCC{8}$; gives a complete axiomatization in hybrid modal logic;
    and identifies the coincidence obstruction with explicit constructions.

  \item \textbf{The even-odd chain:}
    Lutz and Wolter's successor construction (using alternating markers along
    a $tpp$-chain to simulate a functional successor) is essentially the same
    as Wessel's even-odd chain~\cite{Wessel2003}.  Lutz and Wolter credit it
    to Marx--Reynolds~\cite{MarxReynolds1999} and
    Rybina--Zakharyaschev~\cite{RybinaZakharyaschev2001}, who used similar
    techniques in temporal logic.  Wessel's application of the technique to
    RCC8 description logics---which predates Lutz--Wolter by three to four
    years---is not acknowledged.

  \item \textbf{The grid construction and coincidence obstruction:}
    Wessel's Figure~10~\cite{Wessel2003} constructs explicit $n \times n$
    grid models from RCC8 networks and shows they are composition-consistent.
    His Figure~11 identifies the coincidence obstruction: the composition
    table allows multiple relations at the diagonal point, so the grid
    cannot be enforced by a concept term.  These are the same issues Lutz and
    Wolter must address (and solve, but only in the topological setting).
    The prior analysis is not cited.

  \item \textbf{The undecidability boundary:}
    Wessel explicitly characterizes $\ALCIRCC{8}$ as ``close to
    undecidability''~\cite{Wessel2003} and identifies the hybrid logic
    binding operator as the precise extension that crosses the boundary.
    This analysis directly anticipates the structure of the Lutz--Wolter
    proof: they succeed precisely where Wessel showed the abstract semantics
    fails, by exploiting topological rigidity in place of hybrid logic
    binding.
\end{enumerate}

\begin{table}[t]
  \centering
  \caption{Comparison of technical ingredients}
  \label{tab:comparison}
  \smallskip
  \begin{tabular}{@{}lcc@{}}
    \toprule
    \textbf{Construction} & \textbf{Wessel (2002/2003)} & \textbf{Lutz--Wolter (2006)} \\
    \midrule
    Even-odd chain via TPP/NTPP & \cite{Wessel2003}, \S4.4.2 & \cite{LutzWolter2006}, \S8 \\
    Binary counter ($2^n$ path) & \cite{Wessel2003}, \S4.4.2 & --- \\
    $\EXPTIME$-hardness of RCC8 logic & \cite{Wessel2003}, Thm.\ implicit & --- \\
    Grid model construction & \cite{Wessel2003}, Fig.~10 & \cite{LutzWolter2006}, \S8 \\
    Coincidence obstruction identified & \cite{Wessel2003}, Fig.~11 & --- \\
    Coincidence resolved (topologically) & --- & \cite{LutzWolter2006}, Thm.~8.2 \\
    Domino reduction completed & --- & \cite{LutzWolter2006}, Thm.~8.2 \\
    \bottomrule
  \end{tabular}
\end{table}

We emphasize that the Lutz--Wolter contribution is genuine: showing that
topological rigidity resolves the coincidence obstruction is a non-trivial
result.  But the prior identification of the obstruction itself, the
construction of explicit grid models, and the even-odd chain technique applied
to RCC8 logics are Wessel's contributions and should have been cited
accordingly.

%% ====================================================================
\section{Implications for $\ALCIRCC{8}$ decidability}
\label{sec:implications}
%% ====================================================================

The non-transfer result has several implications.

\subsection{The standard reduction is blocked}

The domino tiling reduction---the most common tool for proving description
logic undecidability---cannot be applied to $\ALCIRCC{8}$ via the
Lutz--Wolter route.  The coincidence obstruction blocks it at the completeness
step.

\subsection{Other reductions are also blocked}

More generally, every known undecidability proof for description logics
ultimately encodes a two-dimensional grid.  Grid encoding requires either:
\begin{itemize}
  \item functional roles (not available: the role box is fixed),
  \item number restrictions (not available),
  \item role intersection or role composition beyond the fixed table (not available), or
  \item topological/geometric rigidity (available in $\LRCC{8}$, not in $\ALCIRCC{8}$).
\end{itemize}
As observed in~\cite{OurMainPaper}, the patchwork property of RCC8 actively
\emph{resists} grid encoding: local consistency implies global consistency,
leaving no room for the rigid global constraints that tiling reductions require.

\subsection{The alternating-type trick: what $\ALCIRCC{8}$ CAN do}

While $\ALCIRCC{8}$ cannot encode grids, it is structurally more expressive
than $\ALCIRCC{5}$ in one key respect: the $\forall TPP / \forall NTPP$
distinction provides a limited form of counting along chains.  This is enough
for:
\begin{itemize}
  \item Exponential-length paths (binary counter, $\EXPTIME$-hardness);
  \item A ``1.5-dimensional'' structure: two parallel TPP-chains connected by PO
    rungs (the ``$2 \times \infty$ ladder'' of~\cite{OurMainPaper});
  \item Monotonic broadcast: $\forall NTPP.D$ propagates $D$ to all indirect
    successors.
\end{itemize}
But it is \emph{not} enough for:
\begin{itemize}
  \item Two-dimensional grids (coincidence obstruction);
  \item Consumable communication (monotonicity prevents it);
  \item Unbounded synchronization (required by domino tiling and PCP).
\end{itemize}

This gap---enough for $\EXPTIME$-hardness but not enough for
undecidability---places $\ALCIRCC{8}$ in a genuinely intermediate position.
Whether the logic is decidable remains, after more than twenty years, one of
the most interesting open problems in spatial description logics.

\subsection{Relationship to $\ALCIRCC{5}$}

The non-transfer result for $\ALCIRCC{8}$ is structurally parallel to the
non-transfer of $\LRCC{5}(\RSE)$ undecidability to $\ALCIRCC{5} = \LRCC{5}(\RS)$
(Remark~\ref{rem:rcc5-parallel}).  In both cases, the abstract
composition-table semantics lacks a rigidity property that the concrete
semantics provides:

\begin{center}
\begin{tabular}{@{}lll@{}}
  \toprule
  \textbf{Logic} & \textbf{Concrete rigidity} & \textbf{Missing in abstract} \\
  \midrule
  $\LRCC{8}$ & Topological geometry of $\reals^2$ & Grid coincidence \\
  $\LRCC{5}(\RSE)$ & Supremum closure & Supremum regions \\
  \bottomrule
\end{tabular}
\end{center}

Both $\ALCIRCC{5}$ and $\ALCIRCC{8}$ remain below the known undecidability
boundary.  Whether they are also above the decidability boundary is the
open question.

%% ====================================================================
\section{Conclusion}
%% ====================================================================

We have shown that the Lutz--Wolter undecidability proof for $\LRCC{8}$ does
not transfer to $\ALCIRCC{8}$.  The obstruction is the
\emph{grid coincidence condition}: topological geometry forces the east-of-north
to equal the north-of-east, but the abstract composition table does not.

This obstruction was first identified by Wessel in~\cite{Wessel2003}, who:
\begin{itemize}
  \item constructed explicit RCC8 grid models (Figure~10),
  \item identified the coincidence gap (Figure~11),
  \item developed the even-odd chain for functional successors and the binary
    counter for $\EXPTIME$-hardness (Section~4.4.2),
  \item characterized $\ALCIRCC{8}$ as ``close to undecidability'' (Section~4.4.3),
  \item identified the hybrid logic binding operator as the precise boundary
    (p.~43--44).
\end{itemize}
These contributions predate Lutz--Wolter~\cite{LutzWolter2006} by three to
four years.  While the Lutz--Wolter result is genuinely new (resolving the
coincidence in the topological setting is non-trivial), the prior analysis that
\emph{identified} the obstruction and \emph{demonstrated} that it blocks the
reduction in the abstract setting deserved more explicit
acknowledgment.

The decidability of $\ALCIRCC{8}$ remains open.  The coincidence obstruction
is evidence \emph{for} decidability---it shows that the logic resists the
standard undecidability machinery---but it is not a proof.  Settling the
question will require either a genuinely new undecidability technique that
bypasses the coincidence gap, or a decision procedure that tames the infinite
complete-graph models.

\bigskip

\begin{thebibliography}{99}

\bibitem{LutzWolter2006}
C.~Lutz and F.~Wolter.
\newblock Modal logics of topological relations.
\newblock \emph{Logical Methods in Computer Science}, 2(2):Article~5, 2006.
\newblock DOI: 10.2168/LMCS-2(2:5)2006.

\bibitem{Wessel2001}
M.~Wessel.
\newblock Obstacles on the way to qualitative spatial reasoning with
  description logics: Some undecidability results.
\newblock In \emph{Proc.\ DL Workshop 2001}, CEUR-WS \#49, pp.~96--105, 2001.

\bibitem{Wessel2003}
M.~Wessel.
\newblock Qualitative spatial reasoning with the $\mathcal{ALCI}_{\mathit{RCC}}$
  family -- first results and unanswered questions.
\newblock Technical Report FBI-HH-M-324/03, University of Hamburg, Computer
  Science Department, 2002/2003.

\bibitem{MarxReynolds1999}
M.~Marx and M.~Reynolds.
\newblock Undecidability of compass logic.
\newblock \emph{Journal of Logic and Computation}, 9(6):897--914, 1999.

\bibitem{RybinaZakharyaschev2001}
T.~Rybina and D.~Zakharyaschev.
\newblock Computational complexity of spatial logics.
\newblock Manuscript, 2001.

\bibitem{RenzNebel1999}
J.~Renz and B.~Nebel.
\newblock On the complexity of qualitative spatial reasoning: A maximal
  tractable fragment of the Region Connection Calculus.
\newblock \emph{Artificial Intelligence}, 108(1--2):69--123, 1999.

\bibitem{OurMainPaper}
Claude and M.~Wessel.
\newblock Decidability of $\mathcal{ALCI}_{\mathit{RCC5}}$: An EXPTIME upper
  bound via quasimodel elimination.
\newblock Technical report, 2025.

\end{thebibliography}

\end{document}
